\chapter{Introduction}

Modern machine learning (ML) systems are increasingly embedded in sensitive decision-making pipelines such as healthcare diagnostics, fraud detection, cybersecurity monitoring, and financial risk assessment. As these models grow in influence, the need for verifiable and trustworthy outputs becomes critical. In many practical deployments, a client interacts with a remote model via an API or cloud service without access to the underlying parameters or computation process. This lack of transparency raises concerns regarding trust, fairness, and correctness, especially when decisions directly impact user safety, financial transactions, or regulatory compliance~\cite{lavin2024surveyapplicationszeroknowledgeproofs}.

Zero-Knowledge Proofs (ZKPs) offer a cryptographic solution to these challenges by enabling verifiable computations without revealing internal model details. This thesis explores how ZKPs can be applied to logistic regression to enable privacy-preserving verifiable predictions and accuracy reporting~\cite{cryptoeprint:2025/507}.

\section{Background}

Zero-Knowledge Proofs, formalized by Goldwasser, Micali, and Rackoff~\cite{goldwasser1989knowledge}, are cryptographic protocols that allow a prover to convince a verifier that a statement is correct without revealing any additional information beyond its correctness. A ZKP system satisfies three essential properties:
\begin{itemize}
    \item \textbf{Completeness}: An honest prover always convinces the verifier.
    \item \textbf{Soundness}: A dishonest prover cannot produce a valid proof for a false statement except with negligible probability.
    \item \textbf{Zero-knowledge}: The proof leaks no information about secret values or intermediate computations.
\end{itemize}

These properties make ZKPs ideal for machine learning environments where revealing the model structure, parameters, or training data would be undesirable or unsafe. In particular, ZKPs enable verifiable inference and evaluation without compromising confidentiality~\cite{cryptoeprint:2023/1345}.

\section{Motivation and Problem Statement}

Traditional approaches to verification may require sharing internal model parameters, openly exposing datasets, or depending on a trusted auditor---all of which compromise privacy. ZKPs remove the need for trust by enabling the prover to demonstrate correct computation while keeping all sensitive information hidden~\cite{kates2020plonk}.

The core problem addressed in this report is:
\begin{quote}
\textit{How can Zero-Knowledge Proofs be used to design an efficient, scalable, privacy-preserving framework for logistic regression that allows for verifiable predictions in batch without revealing model parameters?}
\end{quote}

\section{Project Scope and Contributions}

This project is divided into two phases. Phase 1 established the theoretical and practical foundations of verifying a logistic regression model using a PLONK-based arithmetic circuit. Phase 2 focused on scaling the framework for real-world deployment. The key contributions of this work include:
\begin{itemize}
    \item \textbf{Circuit Optimization:} Refined the sigmoid lookup table precision bounds, reducing the constraint footprint by 32\% and cutting compile and setup times nearly in half.
    \item \textbf{Batch Parallel Proving:} Implemented a highly concurrent worker-pool pipeline that evaluates multiple dataset samples in parallel, amortizing setup times and achieving up to a 7.1$\times$ throughput speedup on HPC clusters.
    \item \textbf{N-Feature Generalization:} Restructured the underlying circuit constraints to dynamically scale to $N$-feature models, generalizing the framework beyond specific domain use-cases.
    \item \textbf{Smart Contract Integration Design:} Designed an architecture for exporting the verification key to a Solidity smart contract, enabling trustless, on-chain verification of predictions generated by a private off-chain provider.
\end{itemize}

These advancements demonstrate that cryptographic guarantees of model inference can be practically scaled, paving the way for trustless ML marketplaces and privacy-first prediction systems.
